\documentclass[12pt, titlepage]{article}

\usepackage{amsmath, mathtools}

\usepackage[round]{natbib}
\usepackage{amsfonts}
\usepackage{amssymb}
\usepackage{graphicx}
\usepackage{colortbl}
\usepackage{xr}
\usepackage{hyperref}
\usepackage{longtable}
\usepackage{xfrac}
\usepackage{tabularx}
\usepackage{float}
\usepackage{siunitx}
\usepackage{booktabs}
\usepackage{multirow}
\usepackage[section]{placeins}
\usepackage{caption}
\usepackage{fullpage}

\hypersetup{
bookmarks=true,     % show bookmarks bar?
colorlinks=true,       % false: boxed links; true: colored links
linkcolor=red,          % color of internal links (change box color with linkbordercolor)
citecolor=blue,      % color of links to bibliography
filecolor=magenta,  % color of file links
urlcolor=cyan          % color of external links
}

\usepackage{array}

\externaldocument{../../SRS/SRS}

\input{../../Comments.text}
\input{../../Common.text}

% Module numbers, the same as the MG (Table 1 of the MG). \mref{mAudio} gives
% M3. The labels are set in the Module Decomposition table below.
\newcounter{mnum}
\newcommand{\mthemnum}{M\themnum}
\newcommand{\mref}[1]{M\ref{#1}}
\newcommand{\mdef}[1]{\refstepcounter{mnum}\mthemnum\label{#1}}

% Table of a module's exported access programs. Use one row per program:
% name & inputs & output & exceptions \\
\newenvironment{accessprograms}
  {\par\noindent\small\renewcommand{\arraystretch}{1.2}%
   \tabularx{\textwidth}{@{}>{\raggedright\arraybackslash}p{0.21\textwidth}
   >{\raggedright\arraybackslash}X
   >{\raggedright\arraybackslash}p{0.17\textwidth}
   >{\raggedright\arraybackslash}p{0.24\textwidth}@{}}
   \toprule \textbf{Name} & \textbf{In} & \textbf{Out} & \textbf{Exceptions}\\
   \midrule}
  {\bottomrule\endtabularx\par}

\begin{document}

\title{Module Interface Specification for \progname{}}

\author{\authname}

\date{\today}

\maketitle

\pagenumbering{roman}

\section{Revision History}

\begin{tabularx}{\textwidth}{>{\raggedright\arraybackslash}p{3cm}
>{\raggedright\arraybackslash}p{3.3cm}>{\raggedright\arraybackslash}X}
\toprule \textbf{Date} & \textbf{Developer(s)} & \textbf{Change}\\
\midrule
11 November 2026 & Ahmed Elzaria & Added the introduction, the notation, the
types shared between modules and the module decomposition, and gave each module
its own file (\#324).\\
11 November 2026 & Ahmed Elzaria & Added the MIS of the Audio Output, Lesson
Catalogue and Feedback Modules, and the feedback, image and sound types
(\#325).\\
11 November 2026 & Ahmed Elzaria & Added the MIS of the Attempt, Attempt
Evaluation and Geometry Modules (\#326).\\
11 November 2026 & Ahmed Elzaria & Added the MIS of the Guided Lesson Control
and Improvement Measure Modules, with the access program for each lesson
transition of SRS Section~9.2 (\#327).\\
11 November 2026 & Ahmed Elzaria & Added the MIS of the Attempt Record Module,
and checked the lesson and attempt modules against the MIS checklist
(\#328).\\
\bottomrule
\end{tabularx}

~\newpage

\section{Symbols, Abbreviations and Acronyms}

The symbols, abbreviations and acronyms are those of the SRS~\citep{SRS}
(Section~4, Naming Conventions and Terminology) and of the
MG~\citep{MG} (Section~2, Reference Material). Symbolic constants, such as
\texttt{STAGE\_PASS} and \texttt{ATTEMPT\_LIMIT}, are defined with their units,
ranges and defaults in the Symbolic Constants subsection of the SRS. This
document adds the following.

\begin{center}
\renewcommand{\arraystretch}{1.2}
\begin{tabular}{l p{0.7\textwidth}}
\toprule
\textbf{Symbol} & \textbf{Description}\\
\midrule
M & Module, numbered as in the MG (for example, \mref{mAudio})\\
MIS & Module Interface Specification\\
$D$, $T$, $E$, $I$ & The stages of a guided lesson: demonstration, full-outline
tracing, endpoint-only guidance and independent drawing (SRS
Section~9.2)\\
\bottomrule
\end{tabular}
\end{center}

\newpage

\tableofcontents

\newpage

\pagenumbering{arabic}

\section{Introduction}

The following document details the Module Interface Specifications (MIS) for
\progname{}, a website on which a young child practises drawing lines, shapes
and simple drawings with a stylus. In guided lessons the guidance is removed
stage by stage, and in the creative space the child draws freely and can have a
drawing refined. A supporting adult chooses the lessons, sets the values that
control them and sees the child's progress.

Complementary documents include the System Requirement Specifications
(SRS)~\citep{SRS} and the Module Guide (MG)~\citep{MG}. The MG names each module
and the secret it hides; this document gives each module's interface. The full
documentation and implementation can be found at
\url{https://github.com/AHMEDELZARIA/squigglr}.

\section{Notation}

The structure of the MIS for modules comes from \citet{HoffmanAndStrooper1995},
with the addition that template modules have been adapted from
\cite{GhezziEtAl2003}.  The mathematical notation comes from Chapter 3 of
\citet{HoffmanAndStrooper1995}.  For instance, the symbol := is used for a
multiple assignment statement and conditional rules follow the form $(c_1
\Rightarrow r_1 | c_2 \Rightarrow r_2 | ... | c_n \Rightarrow r_n )$.

The following table summarizes the primitive data types used by \progname.

\begin{center}
\renewcommand{\arraystretch}{1.2}
\noindent
\begin{tabular}{l l p{7.5cm}}
\toprule
\textbf{Data Type} & \textbf{Notation} & \textbf{Description}\\
\midrule
Boolean & $\mathbb{B}$ & true or false\\
character & char & a single symbol or digit\\
integer & $\mathbb{Z}$ & a number without a fractional component in (-$\infty$, $\infty$) \\
natural number & $\mathbb{N}$ & a number without a fractional component in [1, $\infty$) \\
natural number with zero & $\mathbb{N}_0$ & a number without a fractional
component in [0, $\infty$)\\
real & $\mathbb{R}$ & any number in (-$\infty$, $\infty$)\\
\bottomrule
\end{tabular}
\end{center}

\noindent
The specification of \progname \ uses some derived data types: sequences,
strings, sets, tuples and enumerations. Sequences are lists filled with
elements of the same data type, written ``seq of $T$''. Strings are sequences
of characters. Sets are unordered collections without repeats, written ``set
of $T$''. Tuples contain a list of values, potentially of different types, and
their fields are named, as in $p.x$. An enumeration is a type with a fixed list
of values, written as the set of those values, such as $\{D, T, E, I\}$. In
addition, \progname \ uses functions, which are defined by the data types of
their inputs and outputs. Local functions are described by giving their type
signature followed by their specification.

The symbolic constants of the SRS keep their names, written in typewriter font,
for example \texttt{STAGE\_PASS}. A module whose behaviour depends on an
adult-set value (the nine values of BR-008 in the SRS) is given the value as
an input. Only the Guided Lesson Control Module (\mref{mLessonControl}) reads
these values from the Learner Settings Module (\mref{mSettings}), when a lesson
is chosen, and it passes them on to the modules it calls, so that those modules
never read the Learner Settings Module themselves. A constant that the team
fixes belongs to the one module that uses it. It is an exported constant of
that module when its value is known, as \texttt{TRACE\_TOL} is, and part of the
module's secret while the SRS still leaves its value to be decided.

\subsection{Types Shared Between Modules}
\label{sec:shared-types}

Table~\ref{tab:shared-types} lists the types that appear in the interfaces of
more than one module, so that every module uses the same names. They are plain
data types, so using one does not make a module use another module. The last
column names the module whose data the type describes and which decides its
values. StrokeT is the exception: it is an abstract data type, so a module that
calls its access programs uses the Stroke Module (\mref{mStroke}). A module
that only passes a stroke on, as the Guided Lesson Control Module
(\mref{mLessonControl}) does, does not. A type that one module defines and
other modules only name is defined in that module's Exported Types, and the
others refer to it there.

{\small
\renewcommand{\arraystretch}{1.2}
\begin{longtable}{@{}>{\raggedright\arraybackslash}p{0.2\textwidth}
>{\raggedright\arraybackslash}p{0.56\textwidth}
>{\raggedright\arraybackslash}p{0.16\textwidth}@{}}
\caption{Types shared between modules}\label{tab:shared-types}\\
\toprule
\textbf{Type} & \textbf{Definition} & \textbf{Owner}\\
\midrule
\endfirsthead
\toprule
\textbf{Type} & \textbf{Definition} & \textbf{Owner}\\
\midrule
\endhead
\bottomrule
\endlastfoot
PointT & tuple of $(x: \mathbb{R}, y: \mathbb{R})$, a position on the
drawing area & Notation\\
StrokeT & An abstract data type: the points of one stylus contact, in the order
drawn, with the time of each point & \mref{mStroke}\\
PathT & seq of PointT, a line through the points in order. The parts of a
lesson's figure are paths. & Notation\\
DateTimeT & A date and time, in Coordinated Universal Time & Notation\\
StageT & $\{D, T, E, I\}$, the stages of a guided lesson (the set $L$ of SRS
Section~9.2) & \mref{mCatalogue}\\
DrawingStageT & $\{T, E, I\}$, the stages at which the child draws (the set
$A$ of SRS Section~9.2) & \mref{mCatalogue}\\
LessonIdT & Identifies one lesson in the catalogue & \mref{mCatalogue}\\
DifficultyT & tuple of $(\mathit{kind}: \{\mathit{missedPart},
\mathit{missedEndpoint}, \mathit{offFigure}\}, \mathit{index}: \mathbb{N}_0)$,
one predefined difficulty (BR-007): a required part or endpoint of the figure
that was not drawn, given by its index in the figure, or strokes drawn away
from the figure (index 0) & \mref{mCatalogue}\\
FigureT & tuple of $(\mathit{parts}: \text{seq of PathT},
\mathit{endpoints}: \text{seq of PointT})$, the required parts and endpoints of
a lesson's figure (SRS Table~\ref{tab:lesson-catalogue}) & \mref{mCatalogue}\\
GuideT & tuple of $(\mathit{outline}: \text{seq of PathT},
\mathit{endpoints}: \text{seq of PointT})$, what is shown to guide the child at
a drawing stage & \mref{mCatalogue}\\
ClassT & $\{\mathit{complete}, \mathit{incomplete}, \mathit{invalid}\}$ (the
set $R$ of SRS Section~9.2) & \mref{mEvaluation}\\
GuidanceT & $\{\mathit{assisted}, \mathit{independent}\}$ (the set $M$ of SRS
Section~9.2) & \mref{mEvaluation}\\
EvaluationT & tuple of $(\mathit{class}: \text{ClassT}, \mathit{guidance}:
\text{GuidanceT}, \mathit{result}: \mathbb{R}, \mathit{passed}: \mathbb{B},
\mathit{difficulties}: \text{set of DifficultyT})$, the outcome of evaluating
the strokes of an attempt. The result is on the 0--100 scale of the SRS. &
\mref{mEvaluation}\\
ObservedT & tuple of $(\mathit{difficulty}: \text{DifficultyT},
\mathit{corrected}: \mathbb{B})$, a predefined difficulty seen during an
attempt and whether the child corrected it later in the same attempt
(FR-REC-003) & \mref{mAttempt}\\
FeedbackKindT & $\{\mathit{stagePassed},$ $\mathit{lessonComplete},$
$\mathit{keepTrying},$ $\mathit{notEvaluable},$ $\mathit{checkIn},$
$\mathit{lessonEnded}\}$, the kinds of message the child can be given during a
guided lesson, other than the feedback on a predefined difficulty &
\mref{mFeedback}\\
FeedbackT & tuple of $(\mathit{message}: \text{string}, \mathit{image}:
\text{ImageIdT}, \mathit{colour}: \text{ColourIdT}, \mathit{sound}:
\text{SoundIdT})$, the content of one message to the child. The colour is the
one used to show the message and to mark where a missed part or endpoint is &
\mref{mFeedback}\\
ColourIdT & Identifies one colour that \progname{} uses for feedback &
\mref{mFeedback}\\
ImageIdT & Identifies one image that \progname{} shows the child &
\mref{mFeedback}\\
SoundIdT & Identifies one sound that \progname{} can play & \mref{mFeedback}\\
AdultIdT & Identifies a signed-in supporting adult & \mref{mAccess}\\
LearnerIdT & Identifies one learner profile & \mref{mProfile}\\
SettingsT & tuple of the nine adult-set values of one learner (BR-008), each
field named after its SRS constant: \emph{stagePass}, \emph{difficultyStreak},
\emph{attemptLimit}, \emph{improveWindow}, \emph{attemptExtension},
\emph{attemptLimitMax}, \emph{attemptIdle}, \emph{lessonIdlePrompt} and
\emph{lessonIdleEnd} & \mref{mSettings}\\
\end{longtable}
}

\section{Module Decomposition}

The following table is taken directly from the Module Guide document for this
project~\citep{MG}. The module numbers are the same in both documents. The last
column names the issue in which each module's interface is written.

{\small
\renewcommand{\arraystretch}{1.1}
\begin{longtable}{@{}p{0.22\textwidth} l p{0.42\textwidth} l@{}}
\caption{Module Hierarchy}\label{TblMH}\\
\toprule
\textbf{Level 1} & \textbf{No.} & \textbf{Level 2} & \textbf{MIS part}\\
\midrule
\endfirsthead
\toprule
\textbf{Level 1} & \textbf{No.} & \textbf{Level 2} & \textbf{MIS part}\\
\midrule
\endhead
\bottomrule
\endlastfoot
Hardware-Hiding & \mdef{mInputDevice} & Input Device Module & \#265\\
& \mdef{mDisplay} & Display Module & \#265\\
& \mdef{mAudio} & Audio Output Module & \#264\\
\midrule
Behaviour-Hiding & \mdef{mLessonControl} & Guided Lesson Control Module &
\#264\\
& \mdef{mCatalogue} & Lesson Catalogue Module & \#264\\
& \mdef{mAttempt} & Attempt Module & \#264\\
& \mdef{mEvaluation} & Attempt Evaluation Module & \#264\\
& \mdef{mFeedback} & Feedback Module & \#264\\
& \mdef{mImprovement} & Improvement Measure Module & \#264\\
& \mdef{mRecord} & Attempt Record Module & \#264\\
& \mdef{mCreative} & Creative Drawing Module & \#265\\
& \mdef{mRefineControl} & Refinement Control Module & \#265\\
& \mdef{mChildScreen} & Child Screen Module & \#265\\
& \mdef{mAccess} & Adult Access Module & \#266\\
& \mdef{mProfile} & Learner Profile Module & \#266\\
& \mdef{mSettings} & Learner Settings Module & \#266\\
& \mdef{mSummary} & Progress Summary Module & \#266\\
& \mdef{mAdultView} & Adult View Module & \#266\\
\midrule
Software Decision & \mdef{mStroke} & Stroke Module & \#265\\
& \mdef{mGeometry} & Geometry Module & \#264\\
& \mdef{mHistory} & Edit History Module & \#265\\
& \mdef{mRefineEngine} & Refinement Engine Module & \#265\\
& \mdef{mStorage} & Storage Module & \#266\\
& \mdef{mRecovery} & Connection Recovery Module & \#266\\
& \mdef{mLog} & Audit Log Module & \#266\\
\end{longtable}
}

Each module's section below follows the same layout: Module, Uses, Syntax
(exported constants, exported types and access programs), and Semantics (state variables,
environment variables, assumptions, access routine semantics and local
functions). In Design Revision~$-1$ the access routine semantics are left for
Revision~0. The Uses of each module are those of the uses hierarchy in the MG.

\newpage

% One file per module, in module order, so the three MIS parts (#264, #265,
% #266) each edit their own files and do not conflict. The files end in .text,
% not .tex, so the PDF build does not try to compile them on their own.
\input{Modules/M01-InputDevice.text}
\input{Modules/M02-Display.text}
\input{Modules/M03-AudioOutput.text}
\input{Modules/M04-LessonControl.text}
\input{Modules/M05-LessonCatalogue.text}
\input{Modules/M06-Attempt.text}
\input{Modules/M07-AttemptEvaluation.text}
\input{Modules/M08-Feedback.text}
\input{Modules/M09-ImprovementMeasure.text}
\input{Modules/M10-AttemptRecord.text}
\input{Modules/M11-CreativeDrawing.text}
\input{Modules/M12-RefinementControl.text}
\input{Modules/M13-ChildScreen.text}
\input{Modules/M14-AdultAccess.text}
\input{Modules/M15-LearnerProfile.text}
\input{Modules/M16-LearnerSettings.text}
\input{Modules/M17-ProgressSummary.text}
\input{Modules/M18-AdultView.text}
\input{Modules/M19-Stroke.text}
\input{Modules/M20-Geometry.text}
\input{Modules/M21-EditHistory.text}
\input{Modules/M22-RefinementEngine.text}
\input{Modules/M23-Storage.text}
\input{Modules/M24-ConnectionRecovery.text}
\input{Modules/M25-AuditLog.text}

\newpage

\bibliographystyle {plainnat}
\bibliography {../../../refs/References}

\newpage

\newpage{}

\section{Appendix --- Generative AI Use Disclosure}

\wss{ChatGPT (version 5) was used to brainstorm ideas for the template for this disclose section.}

\subsection{AI Tools Used}

\wss{Give the name of the AI tool(s) used and the version.}

\subsection{How and Where Used}

\wss{Explain what the tool(s) were used for.  Examples include: brainstorming,
explaining a technical concept, reviewing the document, generating or modifying
text, generating or modifying diagrams, identifying errors or omissions or
inconsistencies.}

\wss{For each of the uses, identify which parts of the document were affected.}

\wss{You do not need to report every interaction, but you should disclose the 
uses that materially contributed to your work}

\subsection{Verification of AI-Generated Content}

\wss{\begin{itemize}
    \item Describe how the team checked AI-generated or AI-assisted material
    before including it in the document.
    \item In particular, verify factual claims, technical assertions,
    requirements, calculations, references, and other information for which an
    AI system may produce incorrect or fabricated results.
    \item \textbf{The team is responsible for the accuracy, completeness,
consistency, and appropriateness of the submitted document, regardless of
whether material was generated by a human or an AI system.}
\end{itemize}}

\newpage{}

\section{Appendix --- Reflection}

\wss{Not required for CAS 741 projects}

The information in this section will be used to evaluate the team members on the
graduate attribute of Problem Analysis and Design.

\input{../../Reflection.text}

\begin{enumerate}
  \item What went well while writing this deliverable? 
  \item What pain points did you experience during this deliverable, and how
    did you resolve them?
  \item Which of your design decisions stemmed from speaking to your client(s)
  or a proxy (e.g. your peers, stakeholders, potential users)? For those that
  were not, why, and where did they come from?
  \item While creating the design doc, what parts of your other documents (e.g.
  requirements, hazard analysis, etc), it any, needed to be changed, and why?
  \item What are the limitations of your solution?  Put another way, given
  unlimited resources, what could you do to make the project better? (LO\_ProbSolutions)
  \item Give a brief overview of other design solutions you considered.  What
  are the benefits and tradeoffs of those other designs compared with the chosen
  design?  From all the potential options, why did you select the documented design?
  (LO\_Explores)
  \item (After you have implemented another team's module, which means this
  isn't filled in until after the original deadline). What did you learn by
  implementing another team's module? Were all the details you needed in the
  documentation, or did you need to make assumptions, or ask the other team
  questions?  If your team also had another team implement one of your modules,
  what was this experience like?  Are there things in your documentation you
  could have changed to make the process go more smoothly for when an
  ``outsider'' completes some of the implementation?

  Your reflection should be about concepts learned in your SE education, such as modular decomposition, cohesion, coupling, information hiding, etc.  You should critique the quality and utility of the specification you received.

  You should create a PR for your implemented module.  A link to the PR should be part of your reflection.

\end{enumerate}

\end{document}